% ============================================================
% f2012_liaoning_science_q13_three_views.tex
% 2012 辽宁卷（理科）第 13 题：三视图
% 来源 PDF：普通高考/2012/2012辽宁理.pdf
% 原页：第 1 页（A3），印面 751，第 13 题题干下方
% 主视图 4×1，内部距左端 1、3 处有虚线；
% 侧视图 3×1，内部距左端 0.5、2.5 处有虚线；
% 俯视图 4×3 矩形，内含直径 2 的圆。
% ============================================================
\begin{tikzpicture}[>=Stealth, line width=0.7pt, font=\small,
  x=1.5cm, y=1.5cm,
  dim/.style={line width=0.4pt},
  every node/.style={inner sep=1.5pt}]
  % ---------- 主视图 ----------
  \draw (0,4) rectangle (4,5);
  \draw[dashed] (1,4) -- (1,5);
  \draw[dashed] (3,4) -- (3,5);
  % 主视图尺寸
  \draw[dim,<->] (0,3.55) -- (1,3.55) node[midway,below,font=\tiny,fill=white,inner sep=1pt] {$1$};
  \draw[dim,<->] (1,3.55) -- (3,3.55) node[midway,below,font=\tiny,fill=white,inner sep=1pt] {$2$};
  \draw[dim,<->] (3,3.55) -- (4,3.55) node[midway,below,font=\tiny,fill=white,inner sep=1pt] {$1$};

  % ---------- 侧视图 ----------
  \draw (5.3,4) rectangle (8.3,5);
  \draw[dashed] (5.8,4) -- (5.8,5);
  \draw[dashed] (7.8,4) -- (7.8,5);
  \draw[dim,<->] (5.3,3.55) -- (5.8,3.55) node[midway,below,font=\tiny,fill=white,inner sep=1pt] {$0.5$};
  \draw[dim,<->] (5.8,3.55) -- (7.8,3.55) node[midway,below,font=\tiny,fill=white,inner sep=1pt] {$2$};
  \draw[dim,<->] (7.8,3.55) -- (8.3,3.55) node[midway,below,font=\tiny,fill=white,inner sep=1pt] {$0.5$};
  % 侧视图高度
  \draw[dim,<->] (8.55,4) -- (8.55,5) node[midway,right,font=\tiny,fill=white,inner sep=1pt] {$1$};

  % ---------- 俯视图 ----------
  \draw (0,0) rectangle (4,3);
  \draw (2,1.5) circle (1);
\end{tikzpicture}
